% Rewritten 2026-11-20 (user): short, about half a column; replaces the long 2026-10-11 draft.
  \Ac{lisa} will resolve thousands of \acp{gb} whose signals overlap in frequency, so their
  parameters must be inferred jointly, which a \ac{mcmc} global fit does at great cost and must
  repeat for every new stretch of data. We train a conditional flow-matching model once per band
  that returns the joint posterior of the binaries in a frequency window, under the full
  instrument response and broad priors. A flat chart of eight coordinates per source and slots
  ordered in frequency remove the polarisation, phase and labelling copies of each source; a
  Gaussian frame read from the data brings posteriors decades narrower than the prior to widths of
  order one; and the data are encoded once per window, so every draw reuses them. Dividing the
  band where the width of the sources' responses doubles gives five subproblems of fixed shapes, each
  trained on its own device. On the core subproblem,
  \qtyrange{\PriorsNarrowLowMhz}{\PriorsNarrowHighMhz}{\milli\hertz},
  \FlowMatchingCostDraws{} draws take \qty{\FlowMatchingCostNarrow}{\second} on one GPU, against
  \qtyrange{\McmcCostJointLo}{\McmcCostJointHi}{\hour} for a joint \ac{mcmc} chain of the same
  window.
